\documentclass[12pt,a4paper]{article}
\usepackage[utf8]{inputenc}
\usepackage[french]{babel}
\usepackage[T1]{fontenc}
\usepackage{amsmath}
\usepackage{amsfonts}
\usepackage{amssymb}
\usepackage[pdftex]{graphicx}
\usepackage{comment}
\usepackage{siunitx}
\usepackage[colorinlistoftodos,prependcaption,textsize=tiny]{todonotes}
\usepackage[left=2cm,right=2cm,top=2cm,bottom=2cm]{geometry}
\author{Angelo Ortiz}
\title{Rapport du Projet}

\begin{document}
\begin{titlepage}
  \centering
  \includegraphics[width=0.30\textwidth]{logo.jpg}\par\vspace{1cm}
  {\scshape\LARGE Sorbonne Universit\'e\par}
  \vspace{1cm}
  {\scshape\Large 2I013 : Projet (application)\par}
  \vspace{1.5cm}
  {\Large \bfseries Sujet :\par}
  {\huge\bfseries IA Football\par}
  \vspace{2cm}
  {\Large\itshape Fangzhou Ye\par}
  {\Large\itshape Angelo Ortiz\par}
  \vfill
  
  % Bottom of the page
  {\large Licence d'Informatique\par}
  {\large Ann\'ee 2017/2018\par}
\end{titlepage}
 
%\newpage
\tableofcontents
  
\newpage
  
\part*{Introduction}
\section*{Pr\'esentation}
Le projet a consist\'e en le d\'eveloppement d'intelligences artificielles de 
joueurs de football. Le simulateur du jeu \'etant d\'ej\`a fourni, notre 
travail a \'et\'e d'impl\'ementer des strat\'egies de jeu pour mieux r\'eussir 
les matches.
L’objectif de ce travail a \'et\'e d'apprendre \`a bien mener un projet 
collaboratif long sur un nouvel environnement de d\'eveloppement, \`a savoir 
Python.
\section*{Aper\c{c}u}
Voici un aper\c{c}u du jeu...
...
\section*{Arborescence}
Le projet a \'et\'e organis\'e en un module {\bfseries ia} dirigeant le 
comportement de nos joueurs et en divers fichiers de test.
\subsubsection*{ia/tools}
Ce fichier contient une classe repr\'esentant l'\'etat du jeu qui comporte des 
m\'ethodes nous facilitant l'acc\`es aux diverses carat\'eristiques de cet 
\'etat, ainsi que plusieurs fonctions utilitaires.

\subsubsection*{ia/conditions}
Comme le nom l'indique, ce fichier a \'et\'e d\'edi\'e aux diverses conditions 
qui d\'eclenchent les diff\'erents types de comportement.

\subsubsection*{ia/actions}
Ce fichier comporte les actions de base que l'on peut r\'ealiser \`a chaque 
instant du jeu.

\subsubsection*{ia/behaviour}
En ce qui concerne ce fichier, ici se trouvent les actions complexes, i.e.\ 
compos\'ees de plusieurs actions de base. En effet, pour une certaine phase de 
jeu, un joueur peut prendre diff\'erentes d\'ecisions.

\subsubsection*{ia/strategies}
C'est dans ce fichier que les strat\'egies ont \'et\'e r\'eunies. Ces 
strat\'egies dirigent le comportement des joueurs tout au long d'un 
match.

\subsubsection*{ia/gene\_optimisation}
Ce fichier contient les classes n\'ecessaires pour la mise en \oe uvre d'une 
impl\'ementation des algorithmes g\'en\'etiques.

\subsubsection*{ia/gs\_optimisation}
On retrouve ici le code utilis\'e pour suivre la m\'ethode de recherche en 
grille.

\subsubsection*{ia/}
init ? ML\_strats ? DT\_strats ? 
\todo{Il reste encore des fichiers}

\addcontentsline{toc}{part}{Introduction}
\newpage

\part{Architecture logicielle} %et demarche
Dans cette premi\`ere partie, on vous expliquera les choix faits 
concernant l'architecture logicielle.

\newpage

\part{Strat\'egies}
Dans cette deuxi\`eme partie, on vous d\'etaille le comportement de nos 
diff\'erents joueurs.

Compte tenu de la taille du terrain de jeu, le nombre maximum de joueurs par 
\'equipe est limit\'e \`a quatre. Ainsi, il y a trois cat\'egories de matches 
selon le nombre de joueurs par \'equipe : un, deux ou quatre.

\section*{Un joueur}
Pour l'\'equipe compos\'ee uniquement d'un joueur, on a d\'evelopp\'e un seul 
joueur que l'on appelle {\itshape Fonceur}. Au tout d\'ebut du 
projet, il faisait tout au maximum : il doit 
s'approcher de la balle \`a toute vitesse lorsqu'il ne la contr\^ole pas et il 
frappe avec toute sa puissance dans l'autre cas. On a d\'ecid\'e de garder le 
m\^eme nom, m\^eme si son comportement diff\`ere. En effet, \`a pr\'esent il 
avance avec la balle de petites distances et lorsqu'un adversaire lui fait 
opposition il essaie de le dribbler ; il frappe vers la cage adverse s'il se 
trouve dans la surface de r\'eparation. 

\section*{Deux joueurs}
Cette cat\'egorie a \'et\'e celle sur laquelle on a consacr\'e la plupart du 
projet. On compte trois joueurs fonctionnels que l'on appelle 
{\itshape Attaquant}, {\itshape Gardien} et {\itshape D\'efenseur}. 

\subsubsection*{Attaquant}
Il a une vocation plut\^ot offensive. Il se comporte diff\'eremment selon 
s'il contr\^ole la balle ou pas. Dans la premi\`ere situation, il essaie 
syst\'ematiquement d'avancer sur le terrain avec la balle et de faire un tir 
vers la cage d\`es qu'il se trouve dans la surface de r\'eparation adverse. 
Cependant, s'il rencontre un adversaire lui bloquant le chemin, il essaie de 
faire une passe vers son co\'equipier s'il est sans marquage, il dribble le 
joueur sinon. Lorsqu'il se trouve sans le contr\^ole de la balle ...

\subsubsection*{Gardien}
Ce joueur se comporte exactement pareil que l'attaquant lorsqu'il s'agit 
d'attaquer. En revanche, en situation d\'efensive, il reste \`a une 
distance de sa cage lorsque l'\'equipe adverse contr\^ole la balle et essaie 
d'intercepter la balle lorsque le porteur du ballon se trouve dans sa surface 
de r\'eparation. S'il voit que son co\'equipier va bient\^ot contr\^oler la 
balle ou il en a d\'ej\`a le contr\^ole, il monte dans le terrain pour lui 
proposer une possibilit\'e de passe.

\subsection*{D\'efenseur}
Le d\'efenseur n'a par contre qu'une vocation d\'efensive. En effet, s'il a la 
balle, il s'en d\'efait le plus rapidement possible soit \`a l'aide d'une passe 
vers son co\'equipier si d\'emarqu\'e, soit via un d\'egagement le moins 
risqu\'e possible. Quant \`a son effort d\'efensif, il demeure \`a une certaine 
distance radial \`a partir de sa cage et s'approche de la balle en vu d'une 
interception s'il en suffisamment proche.

\section*{Quatre joueurs}
Il s'agit de la cat\'egorie que l'on a trouv\'e la plus int\'eressante en 
raison de la difficult\'e de la gestion des espaces et la r\'epartition des 
fonctions. Tout comme pour la cat\'egorie pr\'ecedente, on compte trois 
joueurs fonctionnels que l'on rep\`ere d'ailleurs par les m\^emes noms. 

\subsubsection*{Attaquant}
Il y a deux diff\'erences en son comportement par rapport \`a son pair 
de la cat\'egorie ci-dessus et celles-ci concernent toutes les deux la 
situation o\`u il se trouve dans sa surface de r\'eparation. La premi\`ere fait 
r\'ef\'erence \`a la r\'ecup\'eration de la balle : il ne se 
pr\'ecipite plus et se comporte tel qu'il le fait dans le milieu du terrain. La 
deuxi\`eme correspond \`a l'effort de marquer les adversaires sans 
marquage.

\subsubsection*{Gardien}
\todo{Apparemment c'est le m\^eme gardien qu'en 2v2}
Ce joueur se comporte exactement pareil que l'attaquant lorsqu'il s'agit 
d'attaquer. En revanche, en situation d\'efensive, il reste \`a une 
distance de sa cage lorsque l'\'equipe adverse contr\^ole la balle et essaie 
d'intercepter la balle lorsque le porteur du ballon se trouve dans sa surface 
de r\'eparation. S'il voit que son co\'equipier va bient\^ot contr\^oler la 
balle ou il en a d\'ej\`a le contr\^ole, il monte dans le terrain pour lui 
proposer une possibilit\'e de passe.

\subsection*{D\'efenseur}
\todo{plus d'interception s'il y a un co\'equipier plus proche de la balle que 
lui (sans la balle) (quid de la sit ou il la controle ?)}
Le d\'efenseur n'a par contre qu'une vocation d\'efensive. En effet, s'il a la 
balle, il s'en d\'efait le plus rapidement possible soit \`a l'aide d'une passe 
vers son co\'equipier si d\'emarqu\'e, soit via un d\'egagement le moins 
risqu\'e possible. Quant \`a son effort d\'efensif, il demeure \`a une certaine 
distance radial \`a partir de sa cage et s'approche de la balle en vu d'une 
interception s'il en suffisamment proche.

\newpage

\part{M\'ethodes d'optimisation}
Dans cette troisi\`eme partie, on vous d\'etaille les diff\'erents algorithmes 
appris tout au long du semestre pour l'am\'elioration des strat\'egies 
initialement propos\'ees.

\section{Recherche en grille}
La premi\`ere m\'ethode d'optimisation pre\'esent\'ee en cours a \'et\'e la 
recherche en grille. Celle-ci est utilis\'ee pour l'optimisation d'une 
action.

Tout d'abord, on doit d\'efinir l'action \`a optimiser et le crit\`ere 
d'optimalit\'e, et rep\'erer les param\`etres associ\'es. Puis on 
obtient un ensemble de valeurs pour chaque param\`etre. 
Pour les param\`etres discrets, il suffit de prendre toutes les valeurs 
possibles, alors que pour les param\`etres continus il faudra faire en amont 
une discr\'etisation, i.e.\ on divise l'intervalle d\'elimit\'e par des bornes 
d\'efinies selon un pas de pr\'ecision. On compte ainsi des vecteurs dont les 
coordonn\'ees correspondent chacune \`a un param\`etre. On fait ensuite une 
recherche exhaustive sur tous les vecteurs obtenus : on teste chaque vecteur 
sous diff\'erents conditions environnementales et on moyenne les \'evaluations 
 du crit\`ere dans lesdites conditions. Finalement, on prend la valeur optimale 
et son vecteur associ\'e.

Il est important de remarquer que la plupart de nos param\`etres sont continus. 
Pour obtenir des valeurs plus pr\'ecises et un meilleur comportement, on a 
besoin d'un pas de discr\'etisation tr\`es petit. Ainsi, on r\'ealise qu'une 
augmentation du nombre de param\`etres \`a optimiser a une tr\`es forte 
influence sur le temps d'ex\'ecution de l'algorithme mis en place. On en 
d\'eduit que le minuscule gain relatif nombre de param\`etres--temps 
d'ex\'ecution ne convient pas au moyen terme.

\section{Algorithmes g\'en\'etiques}
Dans un deuxi\`eme temps, on est pass\'es aux algorithmes g\'en\'etiques.
\todo{TODO}

\section{Arbres de d\'ecision}

\section{Apprentissage automatique}
Dans la derni\`ere partie du semestre, on s'est concentr\'e sur l'apprentissage 
automatique. On nous a pr\'esent\'e les deux grandes classes 
d'algorithmes d'apprentissage : supervis\'es et non supervis\'es, ainsi que 
l'apprentissage par renforcement. En effet, ce dernier type d'apprentissage 
impl\'ement\'e par l'algorithme de \texttt{Q-learning} a \'et\'e mise en place.

\newpage

\part*{Conclusion}
...
\addcontentsline{toc}{part}{Conclusion}

\end{document}
